\documentclass[11pt,a4paper]{article}
\usepackage[utf8]{inputenc}
\usepackage[T1]{fontenc}
\usepackage[polish]{babel}
\usepackage{geometry}
\geometry{margin=1.5cm}
\usepackage{longtable}
\usepackage{makecell}

\begin{document}

\begin{flushright}
Załącznik nr 4a
\end{flushright}

\vspace{-1em}
\begin{center}
\textbf{\large POTWIERDZENIE UZYSKANIA} \\[0.5em]
\textbf{\large EFEKTÓW UCZENIA SIĘ W RAMACH PRAKTYKI ZAWODOWEJ NA PODSTAWIE ZATRUDNIENIA/STAŻU/DZIAŁALNOŚCI GOSPODARCZEJ*}
\end{center}

\vspace{1em}

\noindent Kierunek studiów: Informatyka \\
Student / ka: \dotfill \\
Nr albumu: \dotfill \\
Kierunek studiów: Informatyka \hfill Specjalność: \dotfill \\
\dotfill

\vspace{1em}
\noindent W ramach pracy zawodowej/stażu/działalności gospodarczej* w wymiarze \dots\dots\dots\dots godzin uzyskał/a / nie uzyskał/a * zakładane dla praktyki zawodowej efekty uczenia się:

\vspace{1em}

\renewcommand{\arraystretch}{1.5}
\begin{longtable}{|c|p{10cm}|c|}
\hline
\textbf{Lp.} & \textbf{Efekty uczenia się} & \textbf{Potwierdzenie uzyskania efektów} \\ \hline
\endfirsthead
\hline
\textbf{Lp.} & \textbf{Efekty uczenia się} & \textbf{Potwierdzenie uzyskania efektów} \\ \hline
\endhead

01 & Ma wiedzę na temat sposobu realizacji zadań inżynierskich dotyczących informatyki z zachowaniem standardów i norm technicznych & \makecell{uzyskał/a *\\ uzyskał/a częściowo*\\ nie uzyskał/a*} \\ \hline
02 & Zna technologie, narzędzia, metody, techniki oraz sprzęt stosowane w informatyce & \makecell{uzyskał/a *\\ uzyskał/a częściowo*\\ nie uzyskał/a*} \\ \hline
03 & Zna ekonomiczne, prawne skutki własnych działań podejmowanych w ramach praktyki oraz ograniczenia wynikające z prawa autorskiego i kodeksu pracy & \makecell{uzyskał/a *\\ uzyskał/a częściowo*\\ nie uzyskał/a*} \\ \hline
04 & Zna zasady bezpieczeństwa pracy i ergonomii w zawodzie informatyka & \makecell{uzyskał/a *\\ uzyskał/a częściowo*\\ nie uzyskał/a*} \\ \hline
05 & Pozyskuje informacje odnośnie technologii, metod, technik, sprzętu wymaganego do realizacji powierzonego zadania, posługując się rozmaitymi źródłami literaturowymi i zasobami publikowanymi w języku polskim jak i angielskim & \makecell{uzyskał/a *\\ uzyskał/a częściowo*\\ nie uzyskał/a*} \\ \hline
06 & W oparciu o kontakty ze środowiskiem inżynierskim zakładu, potrafi podnieść swoje kompetencje, wiedzę i umiejętności, co najmniej z dwóch zakresów: zadania dotyczące sprzętu i oprogramowania: np.: programowania, administrowanie siecią komputerową, konserwacja sprzętu i oprogramowania, bieżące usuwanie usterek, administrowanie zasobami informatycznymi, zakładu pracy / instytucji, (e)‑usługami. & \makecell{uzyskał/a *\\ uzyskał/a częściowo*\\ nie uzyskał/a*} \\ \hline
07 & Opracowuje dokumentację dotyczącą realizacji podejmowanych zadań w ramach praktyki, a także referuje ustnie prezentowane w niej zagadnienia & \makecell{uzyskał/a *\\ uzyskał/a częściowo*\\ nie uzyskał/a*} \\ \hline
08 & Potrafi zidentyfikować problem informatyczny występujący w zakładzie pracy / instytucji, opisać go, przedstawić koncepcję rozwiązania i ją zrealizować. & \makecell{uzyskał/a *\\ uzyskał/a częściowo*\\ nie uzyskał/a*} \\ \hline
09 & Potrafi rozwiązać rzeczywiste zadanie inżynierskie z zakresu działalności informatycznej zakładu pracy/instytucji stosując normy i standardy stosowane w informatyce oraz biorąc pod uwagę aspekty środowiskowe i etyczne. & \makecell{uzyskał/a *\\ uzyskał/a częściowo*\\ nie uzyskał/a*} \\ \hline
10 & Pracuje w zespole zajmującym się zawodowo branżą IT & \makecell{uzyskał/a *\\ uzyskał/a częściowo*\\ nie uzyskał/a*} \\ \hline
11 & Przestrzega zasad etyki zawodowej i zgodnie z tymi zasadami korzysta z wiedzy i pomocy doświadczonych kolegów & \makecell{uzyskał/a *\\ uzyskał/a częściowo*\\ nie uzyskał/a*} \\ \hline
12 & Kontaktując się z osobami spoza branży potrafi zarówno pozyskać od nich niezbędne informacje do realizacji planowanego zadania, jak i przekazać im w sposób zrozumiały informacje i opinie z zakresu informatyki & \makecell{uzyskał/a *\\ uzyskał/a częściowo*\\ nie uzyskał/a*} \\ \hline
13 & Dostrzega w praktyce tempo deaktualizacji wiedzy informatycznej oraz skutki działalności informatyków w szczególności ekonomiczne i społeczne & \makecell{uzyskał/a *\\ uzyskał/a częściowo*\\ nie uzyskał/a*} \\ \hline
\end{longtable}

\vspace{1em}
\noindent \textbf{Potwierdzenie merytorycznej oceny dokumentacji złożonej przez studenta ze wskazaniem jakie elementy efektów uczenia się zaliczonych częściowo powinny zostać uzupełnione:}

\vspace{3em}
\noindent\hfill
\begin{minipage}[t]{0.5\textwidth}
\centering
\dotfill \\
\small data, podpis Przewodniczącego Komisji do spraw praktyk
\end{minipage}

\vspace{2em}
\noindent \footnotesize{* niepotrzebne skreślić} \normalsize

\end{document}